\section{Hạn chế và hướng phát triển}\label{sec:limit}
Các kết quả truy hồi dùng caption gắn với một ID ảnh chuẩn, trong khi nhiều ảnh có thể cùng mô tả đúng truy vấn. Nhãn đối tượng và quan hệ từ scene graph không đầy đủ; cách gán truy vấn trực tiếp hoặc gián tiếp cũng phụ thuộc vào chính các nhãn này. Lượt XAI khảo sát 24 caption được chọn theo kết quả của một seed, còn lượt truy vấn trực tiếp có 34 câu đã duyệt; cả hai là mẫu có chủ đích. Khoảng tin cậy được lấy mẫu theo truy vấn để giữ các seed trong cùng cụm, song chưa thay thế được đánh giá trên một mẫu test đại diện hơn. Can thiệp XAI giữ cố định 100 ảnh ứng viên, nên chỉ phản ánh quyết định xếp hạng trong tập đó.

Nguồn hỗ trợ thị giác của từng cạnh chưa được kiểm tra độc lập trên toàn bộ mẫu; lượt xem ảnh mười trường hợp chỉ chấm mức phù hợp của ảnh đứng đầu. Các giải thích chủ yếu gồm cạnh nối đối tượng với ảnh, nên chưa đủ bằng chứng về việc mô hình khai thác nghĩa của quan hệ giữa các đối tượng. Ba seed của QA-Graph thay đổi cả khởi tạo lẫn lấy mẫu gallery; AdaMF-MAT dùng một seed huấn luyện, vì vậy độ bất định của phương pháp này chưa được đo đầy đủ. Bộ 24 truy vấn quan hệ của cấp 2 cũng chưa có rà ảnh thủ công và được đánh giá trong quy trình riêng.

Hướng tiếp theo là xây tập truy vấn gián tiếp có bằng chứng quan hệ được duyệt, kể cả trường hợp ẩn đối tượng để kiểm tra nhu cầu suy luận, rồi so với baseline chỉ dùng đối tượng trên cùng ảnh ứng viên. Với XAI, cần lấy mẫu truy vấn đại diện hơn, kiểm tra trực tiếp nhãn và quan hệ trên ảnh, đồng thời đo ảnh hưởng của can thiệp lên chất lượng truy hồi toàn tập. Các phép thử đó sẽ phân biệt rõ hơn giữa tính trung thực với hàm chấm điểm, độ đúng của bằng chứng và ích lợi thực tế của quan hệ đồ thị.
